% ============================================================================
% GÉNÉRÉ par scripts/exemples-guide.ts — NE PAS ÉDITER À LA MAIN.
% Régénérer : npm run exemples (chaque chiffre est vérifié contre le moteur).
% ============================================================================
\pexemples Trois ménages, du plus simple au plus riche en mécanismes : l'actif seul
(M4), le retraité (M10, montants d'âge et de pension), le couple monoactif
(M7, montant pour conjoint et transfert de crédits). Les crédits valent
\emph{montant} $\times$ \emph{taux du premier palier} (\pc{14} au
Québec, \pc{14.5} au fédéral en 2025).

\medskip\noindent\textbf{M4 --- personne vivant seule, 40~ans, revenu de travail de \D{50000}.}
\emph{Québec.} Revenu imposable : \num{50000} $-$ \num{1420}
(déduction pour travailleur, plafonnée) $-$ \num{465} (RRQ
supplémentaire, déductible) $=$ \num{48115}.
\begin{align*}
\text{impôt brut} &= \pc{14} \times \num{48115} = \num{6736.10} \quad (\text{1\ier{} palier seulement}) \\
\text{montant seul réduit} &= \pos{\num{2128} - \pc{18.75} \times \pos{\num{48115} - \num{42090}}} = \num{998.31} \\
\text{crédits} &= (\num{18571} + \num{998.31}) \times \pc{14} = \num{2739.70} \\
\text{impôt du Québec} &= \num{6736.10} - \num{2739.70} = \num{3996.40}
\end{align*}
\emph{Fédéral.} Revenu imposable : \num{50000} $-$ \num{465}
$=$ \num{49535} (pas de déduction pour travailleur au fédéral).
\begin{align*}
\text{impôt brut} &= \pc{14.5} \times \num{49535} = \num{7182.58} \\
\text{crédits} &= \num{16129} + \num{2511} + \num{247} + \num{655} + \num{1471} = \num{21013} \\
\text{impôt net} &= \num{7182.58} - \num{21013} \times \pc{14.5} = \num{4135.69} \\
\text{après abattement} &= \num{4135.69} \times (1 - \pc{16.5}) = \num{3453.30}
\end{align*}
Les crédits fédéraux : montant personnel de base bonifié (\num{16129}),
cotisations RRQ de base, RQAP et AE (créditables), montant canadien pour
emploi (\num{1471}). L'abattement du Québec retranche
\pc{16.5} de l'impôt fédéral.

\medskip\noindent\textbf{M10 --- retraité vivant seul, 70~ans, revenu de pension privée de \D{20000}.}
\emph{Québec.} Revenu imposable : pension \num{20000} $+$ PSV
imposable \num{8791.14} $=$ \num{28791.14} (le SRG, non imposable, en est
exclu — mais il entre dans le revenu \emph{net} de \num{29891.99} qui
module les crédits).
\begin{align*}
\text{impôt brut} &= \pc{14} \times \num{28791.14} = \num{4030.76} \\
\text{montant combiné} &= \num{2128} + \num{3906} + \num{3470} = \num{9504} \\
&\quad (\text{vivant seul} + \text{âge} + \text{pension; aucune réduction : } \num{29891.99} < \num{42090}) \\
\text{crédits} &= (\num{18571} + \num{9504}) \times \pc{14} = \num{3930.50} \\
\text{impôt du Québec} &= \pos{\num{4030.76} - \num{3930.50}} = \num{100.26}
\end{align*}
\emph{Fédéral.} Impôt brut \num{4174.72} sur \num{28791.14} ; crédits :
BPA \num{16129}, montant en raison de l'âge \num{9028}
(réduit par le revenu net, SRG \emph{inclus} : \num{29891.99}), montant pour
revenu de pension \num{2000}.
\begin{align*}
\text{impôt net} &= \pos{\num{4174.72} - \num{27157} \times \pc{14.5}} = \num{236.95} \\
\text{après abattement} &= \num{236.95} \times (1 - \pc{16.5}) = \num{197.85}
\end{align*}
Impôt fédéral : \D{197.85}.

\medskip\noindent\textbf{M7 --- couple, 45 et 44~ans, revenus de travail de \D{30000} et \D{0}, sans enfant.}
\emph{Québec.} Adulte~1 : revenu imposable \num{28315}, impôt brut
\num{3964.10}, crédits \num{18571} $\times$ \pc{14} $=$ \num{2599.94}.
Adulte~2 (aucun revenu) : impôt brut nul, crédits \emph{inutilisés} de
\num{2599.94} — \textbf{transférés} au conjoint :
\begin{align*}
\text{impôt du Québec} &= \pos{\num{3964.10} - \num{2599.94} - \num{2599.94}} = \num{0}
\end{align*}
\emph{Fédéral.} L'adulte~1 réclame le \textbf{montant pour conjoint} :
$\pos{\num{16129} - \text{revenu net du conjoint } (0)} = \num{16129}$ —
un second montant personnel de base, en pratique.
\begin{align*}
\text{crédits} &= \num{16129} + \num{16129} + \num{1431} + \num{148.20} + \num{393} + \num{1471} = \num{35701.20} \\
\text{impôt net} &= \pos{\num{4311.58} - \num{35701.20} \times \pc{14.5}} = \num{0} \\
\text{après abattement} &= \num{0} \times (1 - \pc{16.5}) = \num{0}
\end{align*}
Impôt fédéral : \D{0} — le couple monoactif ne paie presque rien,
grâce au montant pour conjoint et au transfert québécois des crédits.
